% Nouveau EGA V (cote 25): archive cote 25, site page 144 (PDF page 145).
% typescript. Draft compared: none.
% AI-transcribed by Claude Opus 5.5 (claude-opus-5-5, adaptive thinking, effort high), Message Batches API, batch ega25b (msgbatch_01A7n5YuXubUPNK1zsepi8RG), 2026-11-13; not yet checked by eye.
\scanpage{25/144}
\mspage{\hand{5}}

La colonne provient de la suite exacte de transtitivité pour les
morphismes lisses $X\to Y$ et $Y\to S$ , en \xout{xx} tensorisant par $\underline{O}_{V1}$ (cela reste
exact, car tous les modules de la suite sont loc libres). La ligne
horizontale est un\xout{e} cas particulier d'une suite exacte obtenue chaque
fois que sur X \deleted{\xout{lisse}} sur S, on a une section \unsure{$\phi$} d'un module loc libre F
et qu'on prend le schéma W des zéros; on trouve une suite exacten
\[ \overset{\hand{\vee}}{F}\otimes \underline{O}_X \longrightarrow \underline{\Omega}^1_{X/S}\otimes \underline{O}_W \longrightarrow \underline{\Omega}^1_{W/S} \longrightarrow 0 \]
et \inserted{(si X/S est lisse)} le premier homomorphisme est injectif exactement en les points où
W est lisse sur X avec la ``bonne '' \unsure{dimension} relative (i.e. partout dans
\inserted[typed]{\xout{pratiquement}}
le cas actuel). \xout{Cet homomorphisme} Cette suite exacte est \xout{écrite en}
\xout{par 16 dans le cas où F = $\underline{O}_X$, la généralisation} \xout{(à ceci près} une
conséquence immédiate de la suite exacte
\[ \underline{J}/\underline{J}^2\longrightarrow \underline{\Omega}^1_{X/S}\otimes \underline{O}_W \longrightarrow \underline{\Omega}^1_{X/S} \longrightarrow 0 \]
\margin{16 \textbar} qui figure au par.16 (on y pourrait mettre en corollaire la version
signalée ici). La caractérisation de l'ensemble des points où on
peut mettre un zéro à gauche est contenue dans le critère \over{j}{J}acobien.

Notons que l'on a un isomorphisme canonique
\[ \underline{P}^1_{X/S} = \underline{\Omega}^1_{X/S} + \underline{O}_X \ ,\ \text{d'où}\ \check{\underline{P}}^1_{X/S} = \hand{\unsure{$\mathcal{G}$|$\mathcal{T}$}}_{X/S} + \underline{O}_X \ , \]
d'autre part on \unsure{vérifi\over{e}{é}} que l'homomorphisme \inserted{composé} \hand{$\mu$} du diagramme \hand{\unsure{(1)}} est nul
sur le facteur $\underline{O}_{V1}$ , et sur le facteur $\hand{\unsure{$\mathcal{G}$|$\mathcal{T}$}}_{X/S}\otimes\underline{O}_{V1}$ \inserted{il} se réduit à
l'homomorphisme \inserted{$M(\phi)'$} déduit de la section $M(\phi)$ de $\underline{\mathrm{Sym}}^2(\hand{$\underline{\Omega}$}^1_{X/S})\otimes \underline{O}_{V1}$ déjà
envisagée. Donc en un point x de X, $M(\phi)$ est non dégénéré i.e. $M(\phi)'$
est surjectif sss \hand{$\mu$} est surjectif \inserted{en x}, et on voit sur le diagramme \hand{(1)},
que cela équivaut aussi à dire que \unsure{\hand{$\nu$}} est surjectif \inserted{en x} (car l'un et l'autre
signifie\inserted{nt} que l'homomorphisme canonique de\xout{xx deux} la somme des deux sous
Modules envisagés de $\underline{\Omega}^1_{X/S}\otimes \underline{O}_{V1}$ dans ce derniers est surjectif en
x). On trouve ainsi :

\margin{144}
